\section{作业及答案：结构力学基础与技巧班第10次作业及答案（力法1）}
\begin{analysisbox}
力法先选多余约束，建立基本体系，再由变形协调方程求多余未知力。柔度系数和自由项应由内力图乘或积分得到。
\end{analysisbox}
\subsection*{第 1 页转写}
研究生入学考试辅导丛书结构力学（第二版）

八、练习题

1.基本概念

5-1判断题。力法方程，实际上是力的平衡方程。（：）（福州大学2012）

5-2判断题。图5-125（a）所示桁架结构可选用图5-125（b）所示的体系作为力

法基本体系。（

）（西南交通大学2008）

5-3

图5-126（b）所示结构可作为图5-126（a）所示结构的基本体系。（

（烟台大学2013）

FP

()

(b)

(a)

(b)

图5-125

图5-126

5-4对比图5-127（a）、（b）两个刚架的关系是（：

）。（中南大学2010)

A.内力相同，变形也相同

B.内力相同，变形不同

C.内力不同，变形相同

D.内力不同，变形也不同

5-5

图5-128所示结构的超静定次数为（

）。（中南大学2009）

A. 7

B. 8

C. -9

D. 10

2EI

4EI

EI

EI

2EI

2E1

图5.-127

图5-128

5-6图5-129（b）为图5-129（a）所示结构的力法基本体系，各杆EI=常数，k

为弹簧刚度，则其力法方程的系数、自由项和右端项分别为：811=

$\triangle$IP

。（浙江大学2009）

5-7判断题。图5-130所示结构中，梁AB的截面EI为常数，各链杆的EA相同，

当EI增大时，则梁截面D弯矩代数值Mp增大。（。）（中国矿业大学2011）

5-8图5-131所示两结构中，荷载Fp作用于跨中A、B处，n1、n2均为比例常数

则当ni>nz时，跨中弯矩（

）。（浙江大学2010）

228

\subsection*{第 2 页转写}
FP

FP

2a

(a)

(b)

图5-129

图5-130

B. MA>MB

A. MA=MB

C. MA<MB

D.不确定

5-9结构上没有荷载就没有内力，这个结论在什么情况下适用？在什么情况下不适

用？（四川大学2009）

5-10判断题。图5-132所示两铰拱，已知支座B处下沉了$\triangle$B=0.01f，此时的水平

反力为零，EI常数。（

）（天津大学2011)

5 - 11

)。

图5－133所示结构，为了减小基础所受弯矩，可（

（西南交通大学

2012)

A.减小Ii

B.增大 I2

C.增大I

D.I1、I2同时减小n倍

n,EI

n,EI

'n2EI

n2EI

12

12

图5-131

图5-132

图5-133

2.荷载作用下的计算

5-12用力法计算并作图5-134所示结构的M图。EI=常数。EA=6EI/?。（北京

工业大学2011)

5-13.用力法计算图5-135所示结构，并作其弯矩图，设各杆EI相同。（湖南大学

2012)

5 - 14

用力法作图5-136所示结构M图。（华南理工大学2012)

5 - 15

用力法计算图5-137所示结构，并作M图。各杆EI=常数。（浙江大学2010、

烟台大学2013)

5 - 16

用力法求图5-138所示桁架支座B的反力。各杆EA相同。（山东建筑大学

2008)

5-17

图5-139（a）所示结构，力法计算时采用如图5-139（b）所示的基本结构：

作其M图。（中国矿业大学2010)

\subsection*{第 3 页转写}
研究生入学考试辅导丛书结构力学（第二版）

2kN/m

3kN

EA

10kN

-4kN/m

EI

4m

L1/4

6m

图5-134

图5-135

图5-136

10kN

EI

EI

EI

10kN/m

3m

3m

()

(b)

图5-137

图5-138

图5-139

3.对称性的利用

5-18图5-140所示结构中各杆EI=常数，在给定荷载作用下，HA=

Mo=

（浙江大学2009）

图5-141所示桁架中杆轴力FNl=

5 - 19

，FN2=

（西南交通大学

2010)

5-20

已知图5-142所示对称结构中AD杆A端的弯矩MAD=qL?/36（左侧受拉），

轴力FNAD=一5ql/12，则Mgc=

。（下侧受拉为正）（浙江大学2010）

Fp

10kN

HA

2m

2m

图5-141

图5-140

图5-142

230

\subsection*{第 4 页转写}
5-21判断题。图5-143所示对称结构EI=常数，中点截面C及AB杆内力应满足

M=0，FQ$\ne$0,FN=0,FNAB=0。（

）（西南交通大学2010）

5-22用力法求图5-144所示结构中CD杆的轴力，并作受弯杆件的弯矩图。设受弯

各杆EI=常数，而CD杆的EA=80。（北京交通大学2011)

5-23用力法计算图5-145所示结构，并绘出M图。已知EA=EI/16，EI=常数。

（合肥工业大学2008）

2g

q-16kN/m

EI

EA

EI

EA

2m

图5-143

图5-144

图5-145

5-24

图5-146所示结构，EI=常数，用力法作其M图（利用对称性）。（河海大学

2011)

5 - 25

用力法作图5－147所示结构M图，各杆EI相同。（河海大学2009）

5 - 26

用力法计算图5-148所示结构并作弯矩图。（西安建筑科技大学2010）

10kN/m

EI-常数

56kN

EI

EI

EI

2EI

EI

3m

4m

4m

3m

图5-146

图5-147

图5-148

5-27

用力法计算图5-149所示结构，并作M图。各杆EI=常数，EA=EI/?。（广

州大学2007)

5-28利用结构的对称性作图5-150所示结构的弯矩图。已知EI=常数，链杆EA为

无穷大。（中国矿业大学2009）

5 -29

用力法作图5－151所示结构的弯矩图。EI=常数。（哈尔滨工业大学2013）

5-30

用力法分析图5-152所示刚架，绘M图。设各受弯杆EI值相同。（清华大学

2010)

\subsection*{第 5 页转写}
叫研究生入学考试辅导丛书结构力学（第二版）

EA

C

q勹B

G 1』

20kN

IF

丫G

y

A

7?7?77' 玉

j\_

2

尸

B

＼

L

L

6m



6m

l

1

，

1

,

,

．，

图5-149

图5-150

4. 弹性支承超静定结构的计算

\subsection*{第 6 页转写}
5个、$\times$，实是协调来件

52、V

5-3、$\times$

ABL共线，们可受

因地不可作为美车活构

\subsection*{第 7 页转写}
493

t、hi>hz，假设 h,

MLMB

5-9、静是结构或起静是中助静是部多适用

起静定闷中碍部不用

\subsection*{第 8 页转写}
增大武

，础不

5-12

Mp

41

\subsection*{第 9 页转写}
413.

32

54

[(+*$\times$弹-

:X-3.Vg

13.16

21-72

\subsection*{第 10 页转写}
26-84

23.cb

7.16

\subsection*{第 11 页转写}
40

Mo

Mo

t-16、

小师

40+2/24

(3+>) Fg

(3+22 =0. 854/

\subsection*{第 12 页转写}
[+*+

405

\$X=90

\subsection*{第 13 页转写}
5-8、反对称

Hs= - f

MA

取1-I南离体

Me=0

NA.

好称，取本活构

12

\subsection*{第 14 页转写}
bipzz

\subsection*{第 15 页转写}
410

-1.93

212

31-2

\subsection*{第 16 页转写}
10X 4X280 kwm

361

4805

\subsection*{第 17 页转写}
4807

X2

60+85

5-76.

SIXitDip20X= 12

18

36

18

18

\subsection*{第 18 页转写}
树物半活构

3E1+

3E2

3x=007899

0.07090m

0.0

\subsection*{第 19 页转写}
，，可

40

前[2x4x40x6+bx4x8]---

m图(v-m)

\subsection*{第 20 页转写}
XI=

LOFN

\subsection*{第 21 页转写}
\emph{}

